%% Elements of Nested Fibrational Cosmology
%% Book VI: Effective Continuum and Calculus Legitimacy
%% Final-Pass Canonical Draft
%%
%% Status legend:
%%   [U]  Unconditional  — proved from standing definitions and prior Unc results
%%   [C]  Conditional    — depends on explicitly declared hypotheses
%%   [D]  Definition-structural — structural, no proof obligation
%%   [R]  Remark only    — expository, non-load-bearing
%%   [O]  Open obligation — not yet proved

\documentclass[12pt,oneside]{amsbook}

\usepackage{amsmath, amssymb, amsthm}
\usepackage{mathrsfs}
\usepackage{enumitem}
\usepackage{hyperref}
\usepackage{geometry}
\geometry{margin=1.2in}

%% ---------------------------------------------------------------
%% Theorem environments
%% ---------------------------------------------------------------

\theoremstyle{plain}
\newtheorem{theorem}{Theorem}[section]
\newtheorem{lemma}[theorem]{Lemma}
\newtheorem{proposition}[theorem]{Proposition}
\newtheorem{corollary}[theorem]{Corollary}

\theoremstyle{definition}
\newtheorem{definition}[theorem]{Definition}
\newtheorem{standingrule}[theorem]{Standing Rule}

\theoremstyle{remark}
\newtheorem{remark}[theorem]{Remark}
\newtheorem{scholium}[theorem]{Scholium}

%% Status tag macro
\newcommand{\status}[1]{\textup{[#1]}}

%% ---------------------------------------------------------------
%% Notation
%% ---------------------------------------------------------------
%%  From Books I--V (all freely used):
%%    \Sigma_\infty, O(U), C_{\rho^*}(U), C_n, \mathbf{U}, K^*, T^*
%%    \mathcal{C} (target category), F (realization functor), E (endpoint)
%%
%%  New in Book VI:
%%    \mathcal{F} = \{f_R\}   certified observable family
%%    h_n(U)                  scale-normalized increment at stage n on U
%%    d_n(U)                  discrete gradient datum
%%    J_n(U)                  discrete flux datum
%%    A_n(U)                  discrete accumulator
%%    \phi(U)                 effective smooth representative
%%    \mathcal{I}             continuum interface regime

\begin{document}

\title{Book VI\\[6pt]Effective Continuum and Calculus Legitimacy}
\author{Elements of Nested Fibrational Cosmology}
\date{}
\maketitle

\tableofcontents

\bigskip

%% ---------------------------------------------------------------
\section*{VI.0\quad Purpose}
%% ---------------------------------------------------------------

Book~VI is the interface legitimacy book of the canonical spine.

Books~I--V have established: the observational quotient (Book~I),
local boundary law (Book~II), transport structure (Book~III), the
universal source $\mathbf{U}$ (Book~IV), and the law of lawful descent
to branches (Book~V).  None of these books licenses continuum,
differential, or integral language.  That language must be earned; it
cannot be smuggled in.

The central question Book~VI answers is: \emph{when} is it legitimate
for a branch book to write differential equations, action functionals,
integral operators, smooth fields, or geometric structures?  The
answer is: only when the relevant certified observable families are
asymptotically coherent in the sense made precise here, only on the
certified observable class, and only within a declared continuum
interface regime.

Book~VI establishes this legitimacy through a three-stage progression.
First, it defines certified observable families and the scale-normalized
increments that are their discrete analogues of derivatives.  Second,
it proves that asymptotically coherent observable families admit
effective smooth representatives.  Third, it establishes when
differential and integral notation are legitimate as shorthand for
the asymptotic behavior of the underlying discrete structure.

The governing theorem of the book is the Effective Continuum
Legitimacy Theorem: differential and integral notation are lawful if
and only if they faithfully encode the asymptotic behavior of
scale-normalized increments and accumulators within a declared
continuum interface regime.

Book~VI is a tool-legitimacy book.  It does not prove closure of any
specific endpoint program.  It does not identify which observable
families satisfy the interface hypotheses --- that burden remains
partly branch-dependent.  It licenses the tools; the branches must
verify eligibility.

%% ---------------------------------------------------------------
\section{Imports and Standing Rules}
%% ---------------------------------------------------------------

\textbf{Imports.}  All certified results of Books~I--V.  In
particular: the observational quotient and defect bookkeeping
(Book~I), the boundary law and capacity structure (Book~II), the
Unified Coercive-Transport Inequality and observational-transfer
objects (Book~III), the POT category and universal source $\mathbf{U}$
(Book~IV), and the branch legitimacy law (Book~V).

\begin{standingrule}[\status{D}]\label{sr:continuum-earned}
Continuum language is licensed only as an effective encoding of
certified discrete structure.  It is not primitive.
\end{standingrule}

\begin{standingrule}[\status{D}]\label{sr:no-derivative-without-increment}
Differential notation may not appear in a canonical branch text unless
the relevant scale-normalized increment structure has been certified
in the declared interface regime.
\end{standingrule}

\begin{standingrule}[\status{D}]\label{sr:no-integral-without-accumulation}
Integral notation may not appear unless the relevant certified
accumulation or transport structure has been established.
\end{standingrule}

\begin{standingrule}[\status{D}]\label{sr:interface-locality}
Continuum language is valid only on the certified observable class and
only within the declared continuum interface regime.
\end{standingrule}

\begin{standingrule}[\status{D}]\label{sr:no-asymptotic-promotion}
Refinement suggestiveness is not enough; asymptotic coherence or
effective convergence must be proved at the required level before
continuum language acquires canonical force.
\end{standingrule}

%% ---------------------------------------------------------------
\section{Certified Observable Families and Discrete Structure}
%% ---------------------------------------------------------------

\begin{definition}[\status{D}]\label{def:cert-obs-family}
A \emph{certified observable family} $\mathcal{F} = \{f_R\}_{R \in
\mathrm{Reg}}$ is a collection of observables, one for each admissible
regime $R$, such that: each $f_R \in O_R$ is a realized stabilized
local observable outcome; and the family is transfer-compatible in the
sense of Book~III, Definition~III.2.6.
\end{definition}

\begin{lemma}[\status{U}]\label{lem:cert-fam-quotient-inv}
Certified observable families are invariant under replacement by
quotient-equivalent and branch-equivalent representatives.
\end{lemma}

\begin{proof}
Each $f_R$ is a stabilized local observable outcome, hence a function
on the stabilized quotient $\Sigma_\infty$ (Book~I, Theorem~I.6.1).
Quotient-equivalent representatives are in the same class; therefore
$f_R$ assigns them the same value.  Branch-equivalence is preserved by
the transfer compatibility condition of the family.
\end{proof}

\begin{definition}[\status{D}]\label{def:scale-norm-increment}
A \emph{scale-normalized increment} of a certified observable family
$\mathcal{F}$ at stage $n$ on a region $U$ is
\[
  h_n(U) \;:=\; \frac{f_{R_{n+1}}(U) - f_{R_n}(U)}{\ell_n},
\]
where $\ell_n > 0$ is the declared comparison length at stage $n$.
The scale-normalized increment measures the change in the observable
per unit comparison length under one admissible enlargement step.
\end{definition}

\begin{lemma}[\status{U}]\label{lem:increment-welldef}
Scale-normalized increments, when defined, depend only on certified
comparison data.
\end{lemma}

\begin{proof}
The numerator $f_{R_{n+1}}(U) - f_{R_n}(U)$ involves only realized
stabilized local outcomes, which are quotient-visible by
Lemma~\ref{lem:cert-fam-quotient-inv}.  The denominator $\ell_n$ is
a declared comparison length, hence part of the certified data of the
interface regime.  The quotient depends only on these certified
quantities.
\end{proof}

\begin{definition}[\status{D}]\label{def:discrete-gradient}
The \emph{discrete gradient datum} $d_n(U)$ of a certified observable
family on a region $U$ at stage $n$ is the vector of scale-normalized
increments in all declared comparison directions.
\end{definition}

\begin{definition}[\status{D}]\label{def:discrete-flux}
The \emph{discrete flux datum} $J_n(U)$ on a region $U$ at stage $n$
is the certified flow of observable content across the collar
$\partial_{\rho^*} U$, determined by the LS-2 boundary-law structure
of Book~II.
\end{definition}

\begin{definition}[\status{D}]\label{def:discrete-accumulator}
A \emph{discrete accumulator} $A_n(U)$ on a region $U$ is the finite
sum
\[
  A_n(U) \;:=\; \sum_{k=0}^{n} J_k(U) \cdot \ell_k,
\]
formed by accumulating flux data with declared comparison lengths.
It is the discrete analogue of a line or volume integral.
\end{definition}

\begin{lemma}[\status{U}]\label{lem:flux-pres-ind}
Discrete gradient data and discrete flux data are
presentation-independent once defined on a certified observable
family.
\end{lemma}

\begin{proof}
Both are defined in terms of stabilized local observable outcomes and
the LS-2 boundary law, all of which are quotient-visible.  Therefore
they are independent of the particular presentation representative
chosen within each equivalence class.
\end{proof}

\begin{lemma}[\status{U}]\label{lem:accumulator-inv}
Discrete accumulators formed from certified flux data are invariant
under lawful decomposition refinement.
\end{lemma}

\begin{proof}
Decomposition refinement is an admissible enlargement.  Under lawful
admissible enlargement, the LS-2 boundary law (Book~II,
Theorem~II.4.1) and the defect-ledger additivity (Book~I,
Proposition~I.5.4) together ensure that the flux data and accumulated
totals are preserved: summing over a refined decomposition gives the
same total as summing over the coarser one.
\end{proof}

\begin{definition}[\status{D}]\label{def:refinement-consistent}
A certified observable family is \emph{refinement-consistent} if the
scale-normalized increments $h_n(U)$ stabilize under admissible
refinement: there exists a stage $N$ beyond which $h_n(U)$ changes by
less than any declared tolerance $\varepsilon > 0$ at each further
admissible refinement step.
\end{definition}

\begin{lemma}[\status{U}]\label{lem:refinement-consistent-stable}
Refinement consistency of a certified observable family is preserved
under lawful source-descended comparison maps.
\end{lemma}

\begin{proof}
Under a lawful source-descended comparison map (a transfer map in the
sense of Book~III), the certified observable content is preserved up to
the declared transfer structure.  Therefore if the increments stabilize
in the source regime, the transferred increments stabilize in the target
regime by the transfer-compatibility conditions of Book~III,
Definition~III.2.6.
\end{proof}

%% ---------------------------------------------------------------
\section{Asymptotic Coherence and Effective Smooth Representatives}
%% ---------------------------------------------------------------

\begin{definition}[\status{D}]\label{def:asymptotic-coherence}
A certified observable family $\mathcal{F}$ is
\emph{asymptotically coherent} if the scale-normalized increments
$h_n(U)$ form a Cauchy sequence in the declared comparison norm as $n
\to \infty$: for every declared tolerance $\varepsilon > 0$, there
exists $N$ such that $|h_m(U) - h_n(U)| < \varepsilon$ for all $m, n
\geq N$.
\end{definition}

\begin{definition}[\status{D}]\label{def:effective-smooth}
An \emph{effective smooth representative} of an asymptotically
coherent certified observable family $\mathcal{F}$ is a smooth
function $\phi: \mathcal{R} \to \mathbb{R}$ defined on a declared
interface regime $\mathcal{R}$ such that:
\begin{enumerate}[label=\textup{(\arabic*)}]
  \item the values $\phi(x)$ approximate the certified observable
        data $f_{R_n}$ within the declared tolerance at each stage
        $n$; and
  \item the differential expressions of $\phi$ approximate the
        scale-normalized increments $h_n$ in the declared comparison
        norm.
\end{enumerate}
The effective smooth representative encodes the observable data; it
does not replace or displace it.
\end{definition}

\begin{definition}[\status{D}]\label{def:continuum-interface-regime}
A \emph{continuum interface regime} is a declared pair
$(\mathcal{R}, \mathcal{F})$ consisting of a smooth domain
$\mathcal{R}$ and a certified observable family $\mathcal{F}$ such
that:
\begin{enumerate}[label=\textup{(\arabic*)}]
  \item $\mathcal{F}$ is asymptotically coherent on $\mathcal{R}$;
  \item the scale-normalized increments converge to the relevant
        differential expressions of the effective smooth representative
        on $\mathcal{R}$; and
  \item the discrete accumulators converge to the corresponding
        integral expressions on $\mathcal{R}$.
\end{enumerate}
\end{definition}

\begin{lemma}[\status{U}]\label{lem:incoherence-blocks}
If asymptotic coherence fails, then no effective smooth representative
is canonically licensed.
\end{lemma}

\begin{proof}
An effective smooth representative requires (Definition~\ref{def:effective-smooth})
that differential expressions approximate scale-normalized increments.
If asymptotic coherence fails, the scale-normalized increments do not
converge; therefore no smooth function's derivatives can approximate
them, and no effective smooth representative exists in the declared
regime.
\end{proof}

\begin{theorem}[\status{C}]\label{thm:refinement-coherence}
\textup{(Refinement Coherence Theorem.)}  A certified observable
family satisfying the declared normalization and persistence
conditions is asymptotically coherent under admissible refinement.
\end{theorem}

\begin{proof}
By Definition~\ref{def:asymptotic-coherence}, asymptotic coherence
requires the scale-normalized increments to form a Cauchy sequence.
Under the declared normalization conditions (specifically, that $\ell_n
\to 0$ at a controlled rate relative to the refinement steps) and the
persistence conditions (that the transfer-compatibility of
Lemma~\ref{lem:refinement-consistent-stable} holds), the increments
$h_n(U)$ form a bounded sequence that changes by less than $\varepsilon$
per step once sufficiently refined.  This Cauchy property follows from
the Finite Stabilization Theorem (Book~I, Theorem~I.6.1) applied to the
quotient of increment values: the quotient stabilizes, so the increments
eventually lie in a common equivalence class, forcing convergence.
\end{proof}

\begin{theorem}[\status{C}]\label{thm:effective-rep-existence}
\textup{(Effective Representative Existence Theorem.)}  If a certified
observable family is asymptotically coherent and satisfies the declared
compactness or stabilization conditions on the continuum interface
regime, then it admits an effective smooth representative.
\end{theorem}

\begin{proof}
By asymptotic coherence (Definition~\ref{def:asymptotic-coherence}),
the scale-normalized increments converge in the declared norm.  The
declared compactness or stabilization conditions ensure that this
convergence is uniform on the interface regime.  By a standard
completion argument in the relevant function space, the limiting object
is a smooth function $\phi$ satisfying the approximation conditions of
Definition~\ref{def:effective-smooth}.
\end{proof}

\begin{theorem}[\status{C}]\label{thm:effective-rep-uniqueness}
\textup{(Effective Representative Uniqueness Theorem.)}  Any two
effective smooth representatives of the same certified observable
family are equivalent with respect to the preserved asymptotic
content.
\end{theorem}

\begin{proof}
Suppose $\phi$ and $\phi'$ are two effective smooth representatives of
the same certified observable family on the same interface regime.  By
Definition~\ref{def:effective-smooth}, both $\phi$ and $\phi'$ have
their differential expressions converging to the same scale-normalized
increments $h_n(U)$.  Therefore $\phi$ and $\phi'$ have equal
derivatives in the limit, and since both are smooth on a connected
domain, they agree up to a constant.  The constant is fixed by the
constraint that both approximate the same certified observable values
$f_{R_n}$.  Hence $\phi = \phi'$.
\end{proof}

%% ---------------------------------------------------------------
\section{Differential and Integral Legitimacy}
%% ---------------------------------------------------------------

\begin{proposition}[\status{C}]\label{prop:differential-shorthand}
Whenever an effective smooth representative $\phi$ exists, its
admissible differential expressions must be interpretable as shorthand
for certified scale-normalized increment behavior.
\end{proposition}

\begin{proof}
By Definition~\ref{def:effective-smooth}\textup{(2)}, the differential
expressions of $\phi$ approximate the scale-normalized increments $h_n$.
Therefore any differential expression of $\phi$ is, to within the
declared tolerance, a re-encoding of $h_n$.  It cannot carry more
observational content than the underlying scale-normalized increment.
\end{proof}

\begin{theorem}[\status{C}]\label{thm:differential-legitimacy}
\textup{(Effective Differential Legitimacy Theorem.)}  For a certified
observable family in a lawful continuum interface regime, differential
notation is legitimate if and only if it faithfully encodes the
asymptotic behavior of scale-normalized increments.
\end{theorem}

\begin{proof}
\emph{Sufficiency:} If differential notation faithfully encodes the
asymptotic behavior of scale-normalized increments, then by
Theorem~\ref{thm:effective-rep-existence} an effective smooth
representative exists, and its differential expressions are the
legitimate continuum encoding of the discrete increment data.

\emph{Necessity:} If differential notation does not faithfully encode
the asymptotic behavior of scale-normalized increments, then the
smooth function whose derivatives are that notation is not an effective
smooth representative (it fails condition~(2) of
Definition~\ref{def:effective-smooth}).  Therefore the differential
notation lacks the asymptotic foundation required for legitimacy.
\end{proof}

\begin{corollary}[\status{C}]\label{cor:no-derivative-outside-regime}
No derivative claim has canonical force outside the certified
observable class and declared interface regime.
\end{corollary}

\begin{corollary}[\status{C}]\label{cor:PDE-backed-by-increments}
Any downstream PDE expression has theorem-level canonical force only
to the extent that its differential terms are backed by certified
normalized increment structure.
\end{corollary}

\begin{remark}[\status{R}]\label{rmk:PDE-citation}
Corollary~\ref{cor:PDE-backed-by-increments} is the standard citation
for branch papers that use PDE language.  When a branch book writes a
partial differential equation --- whether Navier--Stokes, Yang--Mills,
Einstein field equations, or any other --- it must cite the continuum
interface regime in which that equation's differential terms are backed
by certified discrete increment structure.  Without such a citation,
the PDE notation has no canonical force within the NFC program.
\end{remark}

\begin{proposition}[\status{C}]\label{prop:integral-shorthand}
Whenever an effective smooth representative $\phi$ exists, its
admissible integral expressions must be interpretable as shorthand for
certified discrete accumulation or transport structure.
\end{proposition}

\begin{proof}
By Definition~\ref{def:effective-smooth}\textup{(1)}, the values of
$\phi$ approximate the certified observable data at each stage.
By Lemma~\ref{lem:accumulator-inv}, discrete accumulators $A_n(U)$
are refinement-invariant.  Integral expressions of $\phi$ therefore
approximate $A_n(U)$ and re-encode discrete accumulation data
without carrying more observational content than the underlying
discrete structure.
\end{proof}

\begin{theorem}[\status{C}]\label{thm:integral-legitimacy}
\textup{(Effective Integral Legitimacy Theorem.)}  For a certified
observable family in a lawful continuum interface regime, integral
notation is legitimate if and only if it faithfully encodes certified
accumulation or transport data.
\end{theorem}

\begin{proof}
Parallel to the proof of Theorem~\ref{thm:differential-legitimacy},
replacing scale-normalized increments with discrete accumulators and
derivatives with integrals throughout.
\end{proof}

\begin{corollary}[\status{C}]\label{cor:no-integral-outside-regime}
No integral identity has canonical force outside the certified
observable class and declared interface regime.
\end{corollary}

\begin{corollary}[\status{C}]\label{cor:variational-from-accumulation}
Variational or conservation-style integral statements are lawful only
insofar as they descend from certified discrete accumulation structure.
This governs all action principles and balance laws in branch books.
\end{corollary}

%% ---------------------------------------------------------------
\section{The Continuum Bridge Schema}
%% ---------------------------------------------------------------

The previous sections establish when continuum notation is legitimate
for a single certified observable family.  Branch books typically need
to \emph{promote} a result from the certified discrete regime to a
continuum statement.  This promotion requires an explicit bridge
schema.

\begin{definition}[\status{D}]\label{def:scoped-certificate}
A \emph{scoped certificate} for continuum promotion is a certified
result whose force is limited to the explicit regime named in its
statement: the specific certified observable family, the specific
interface regime, and the specific tolerance bounds declared at the
time of the result.
\end{definition}

\begin{definition}[\status{D}]\label{def:continuum-bridge}
A \emph{continuum bridge} is a separately proved theorem carrying a
scoped certificate from a declared discrete or pre-continuum regime
into a target continuum regime, with explicit specification of:
\begin{enumerate}[label=\textup{(\arabic*)}]
  \item the source-side certified discrete burden,
  \item the target continuum criterion being promoted to, and
  \item the scope of the promoted claim.
\end{enumerate}
\end{definition}

\begin{theorem}[\status{U}]\label{thm:continuum-bridge-schema}
\textup{(Continuum Bridge Schema Theorem.)}  A branch may promote a
certified discrete result to continuum force only when:
\begin{enumerate}[label=\textup{(\arabic*)}]
  \item a scoped certificate establishes the result on the source
        discrete regime, and
  \item a separately proved continuum bridge carries that certified
        content to the declared continuum target.
\end{enumerate}
Neither ingredient alone suffices.  A scoped certificate alone
establishes only the discrete claim; a continuum bridge alone has no
certified source content to carry.
\end{theorem}

\begin{proof}
A scoped certificate proves only the claim on its declared discrete
regime; it carries no automatic asymptotic or continuum extension
(Standing Rule~\ref{sr:no-asymptotic-promotion}).  A continuum bridge,
by Definition~\ref{def:continuum-bridge}, must be ``separately proved''
--- it cannot be the same proof as the scoped certificate.  Therefore
both are needed.  The combination provides the source certification and
the bridge, jointly licensing the promoted claim.
\end{proof}

\begin{corollary}[\status{U}]\label{cor:finite-no-self-promote}
Finite certification does not self-promote to continuum force.
\end{corollary}

\begin{corollary}[\status{U}]\label{cor:discrete-source-no-endpoint}
Source-side canonical structure does not self-promote to physical
endpoint identification.
\end{corollary}

\begin{corollary}[\status{U}]\label{cor:conditional-no-unconditional}
Branch-conditional bridge success does not self-promote to
unconditional continuum closure.
\end{corollary}

\begin{theorem}[\status{U}]\label{thm:no-smuggling}
\textup{(Continuum No-Smuggling Corollary.)}  Any use of differential,
integral, field, manifold, or variational language beyond the scope of
a certified continuum interface regime is non-canonical and has no
theorem-level force in the NFC program.
\end{theorem}

\begin{proof}
By Theorem~\ref{thm:differential-legitimacy}, differential notation is
legitimate only within a certified interface regime.  By
Theorem~\ref{thm:integral-legitimacy}, integral notation is legitimate
only within a certified interface regime.  By
Theorem~\ref{thm:continuum-bridge-schema}, promotion from discrete to
continuum requires both a scoped certificate and a continuum bridge.
Any use of such notation outside these conditions lacks the legitimacy
foundation and therefore has no theorem-level force.
\end{proof}

\begin{corollary}[\status{C}]\label{cor:heuristic-marking}
A downstream branch may use suggestive continuum language before full
interface certification, but such language must be explicitly marked as
heuristic.  Heuristic use does not constitute a canonical claim.
\end{corollary}

\begin{proof}
Theorem~\ref{thm:no-smuggling} prohibits canonical use of continuum
language outside certified regimes.  However, it does not prohibit
exploratory or motivational language, provided such language is
explicitly labeled as heuristic and not cited as having theorem-level
force.  The marking requirement follows from Book~VII's status
discipline, which prohibits silent status inflation.
\end{proof}

%% ---------------------------------------------------------------
\section{Governing Theorem of Book VI}
%% ---------------------------------------------------------------

\begin{theorem}[\status{C}]\label{thm:governing}
\textup{(Effective Continuum Legitimacy Theorem.)}  For certified
stabilized observables in a lawful continuum interface regime,
differential and integral notation are lawful if and only if they
faithfully encode the asymptotic behavior of scale-normalized
increments and certified accumulators, respectively, within that
declared regime.
\end{theorem}

\begin{proof}
The differential direction is Theorem~\ref{thm:differential-legitimacy}.
The integral direction is Theorem~\ref{thm:integral-legitimacy}.
Together they give the full biconditional over both differential and
integral language.
\end{proof}

\begin{corollary}[\status{C}]\label{cor:joint-legitimacy}
For a certified observable family in a lawful continuum interface
regime, the full calculus of differential and integral operators is
jointly legitimate as asymptotic encoding of the certified discrete
dynamics.
\end{corollary}

\begin{remark}[\status{R}]
\textup{(E-CONT bridge obligations and discrete Sobolev surrogate.)}
\textit{(Pre-canonical remarks have no canonical force;
see Book~VII, Standing Rule~\ref{sr:precanonical-force}.)}  
The pre-canonical program provides two confirmations that
Book~VI's Continuum Bridge Schema (Theorem~\ref{thm:continuum-bridge-schema})
correctly gates the YM and gravity continuum programs:
\begin{enumerate}[label=\textup{(\roman*)}]
  \item \textbf{E-CONT bridge obligations} (v7.35 \S\S11051, 11067):
    the encoding-map bridge ($\Gamma:\mathcal{K}(G_n)\to M$) and
    curvature-proxy bridge (E3-proxy) that appear in the YM continuum
    program (R3.6) and the gravity continuum program (R1.1--R1.2)
    are both instances of the Continuum Bridge Schema.  Book~VI's
    \texttt{thm:continuum-bridge-schema} is therefore the canonical
    formulation that both programs must cite.
  \item \textbf{NS-4B Sobolev embedding surrogate} (v7.35 \S9333):
    The inequality $\|f\|_\infty \leq C_S(s)\|f\|_{H^s}$ with
    table-computable constant $C_S(s) = F_S(s, d_{\mathrm{eff}},
    C_{\mathrm{HK}})$ is derived from the ACA Route~B
    $\to$ Hypercontractivity $\to$ Nash chain from finite NF$_2$
    data alone.  It is the discrete-native surrogate for the
    Sobolev embedding used in the NS and YM continuum arguments,
    and belongs under Book~VI's licensed continuum bridge
    infrastructure.
\end{enumerate}
Both items are confirmed instances of existing Book~VI machinery;
they require no new canonical material.
\end{remark}

%% ---------------------------------------------------------------
\section{Boundary of Book VI}
%% ---------------------------------------------------------------

\textbf{What Book~VI proves.}
\begin{enumerate}
  \item Scale-normalized increments and discrete accumulators
        are well-defined, POT-native, and presentation-independent.
  \item Asymptotically coherent certified observable families
        admit effective smooth representatives, unique up to
        preservation of asymptotic content.
  \item Differential notation is legitimate iff it faithfully
        encodes scale-normalized increment asymptotics within a
        declared interface regime.
  \item Integral notation is legitimate iff it faithfully
        encodes certified accumulator asymptotics within a
        declared interface regime.
  \item The Continuum Bridge Schema governs all promotion from
        discrete to continuum force.
  \item No continuum language is canonical outside a declared
        and certified interface regime.
\end{enumerate}

\textbf{What Book~VI does not prove, and may not be assumed here.}
\begin{enumerate}
  \item \emph{Which observable families satisfy the interface
        hypotheses.}  The burden of verifying asymptotic coherence
        and the compactness or stabilization conditions is branch-specific.
        Branch books must establish these conditions for their own
        observable families before citing the Effective Continuum
        Legitimacy Theorem.
  \item \emph{Branch closures.}  Licensing continuum language does
        not constitute closing any specific branch endpoint program.
        Yang--Mills, Navier--Stokes, and gravity closure still require
        branch-specific bridge theorems.
  \item \emph{Primitive ontology.}  The existence of an effective
        smooth representative is never proof that the continuum is
        ontologically primitive.  The continuum is an effective
        encoding, not a foundation.
  \item \emph{Action principles, gauge theories, or field equations
        by default.}  No specific PDE, Lagrangian, or field equation
        is licensed by this book.  Each requires its own declared
        interface regime and its own bridge theorem.
\end{enumerate}

%% ---------------------------------------------------------------
\section{Handoff to Book VII}
%% ---------------------------------------------------------------

Book~VI establishes when continuum tools are licensed.  Book~VII
establishes how results using those tools may be cited, scoped, and
combined.  The critical connection is that Book~VII's promotion law
governs the same promotion steps that Book~VI's Continuum Bridge Schema
describes structurally: Book~VI tells you when the bridge is possible;
Book~VII tells you how to state and cite the result once it has been
built.

Every branch book that uses results from Book~VI must also conform to
Book~VII's closure schema.  The two books work in tandem: interface
legitimacy and closure discipline are complementary, not alternative,
requirements.

%% ---------------------------------------------------------------
\section{Dependency Ledger}
%% ---------------------------------------------------------------

\renewcommand{\arraystretch}{1.4}
\noindent
\begin{tabular}{@{}p{4.2cm}p{0.8cm}p{3.4cm}p{4.6cm}@{}}
\hline
\textbf{Theorem} & \textbf{St.} & \textbf{Depends on} & \textbf{Exports to} \\
\hline
\ref{lem:cert-fam-quotient-inv}\ Cert.\ fam.\ quotient-inv.
  & [U] & Bk.~I Thm., Bk.~III Def.
  & \ref{lem:increment-welldef}, \ref{thm:governing} \\

\ref{lem:increment-welldef}\ Increments cert.-data
  & [U] & \ref{lem:cert-fam-quotient-inv}, Def.~\ref{def:scale-norm-increment}
  & \ref{thm:differential-legitimacy}, \ref{thm:governing} \\

\ref{lem:flux-pres-ind}\ Flux pres.-indep.
  & [U] & Bk.~II LS-2, Bk.~I quotient
  & \ref{lem:accumulator-inv}, \ref{thm:integral-legitimacy} \\

\ref{lem:accumulator-inv}\ Accumulator inv.
  & [U] & \ref{lem:flux-pres-ind}, Bk.~I ledger add.
  & \ref{thm:integral-legitimacy}, \ref{thm:governing} \\

\ref{lem:incoherence-blocks}\ Incoherence blocks smooth rep.
  & [U] & Defs.~\ref{def:asymptotic-coherence}, \ref{def:effective-smooth}
  & \ref{thm:effective-rep-existence} \\

\ref{thm:refinement-coherence}\ Refinement Coherence Thm.
  & [C] & Defs.~\ref{def:cert-obs-family}--\ref{def:asymptotic-coherence}
  & \ref{thm:effective-rep-existence}, \ref{thm:governing} \\

\ref{thm:effective-rep-existence}\ Eff.\ Rep.\ Existence Thm.
  & [C] & \ref{thm:refinement-coherence}, compactness
  & \ref{thm:effective-rep-uniqueness}, \ref{thm:differential-legitimacy},
    \ref{thm:integral-legitimacy} \\

\ref{thm:effective-rep-uniqueness}\ Eff.\ Rep.\ Uniqueness Thm.
  & [C] & \ref{thm:effective-rep-existence}
  & \ref{thm:governing} \\

\ref{thm:differential-legitimacy}\ Diff.\ Legitimacy Thm.
  & [C] & \ref{thm:effective-rep-existence}, SR~\ref{sr:no-derivative-without-increment}
  & \ref{cor:PDE-backed-by-increments}, \ref{thm:governing} \\

\ref{thm:integral-legitimacy}\ Int.\ Legitimacy Thm.
  & [C] & \ref{lem:accumulator-inv}, \ref{thm:effective-rep-existence}
  & \ref{cor:variational-from-accumulation}, \ref{thm:governing} \\

\ref{thm:continuum-bridge-schema}\ Continuum Bridge Schema Thm.
  & [U] & Def.~\ref{def:continuum-bridge}, SR~\ref{sr:no-asymptotic-promotion}
  & \ref{thm:no-smuggling}, \ref{cor:finite-no-self-promote}; Bk.~VII;
    all branch books \\

\ref{thm:no-smuggling}\ Continuum No-Smuggling Corollary
  & [U] & \ref{thm:differential-legitimacy}, \ref{thm:integral-legitimacy},
    \ref{thm:continuum-bridge-schema}
  & Constitutional rule for all branch books \\

\ref{thm:governing}\ Effective Continuum Legitimacy Thm.
  & [C] & \ref{thm:differential-legitimacy}, \ref{thm:integral-legitimacy}
  & Book~VII; all branch books using calc./PDE language \\
\hline
\end{tabular}
\renewcommand{\arraystretch}{1}

\medskip

\begin{scholium}[\status{R}]\label{sch:book-vi-honest-posture}
Book~VI is the thinnest book in the spine from a source-content
perspective.  Its theorems are mostly biconditionals: differential
notation is legitimate iff it encodes the right thing; integral
notation is legitimate iff it encodes the right thing.  The conceptual
weight lies in the definitions --- what ``encoding'' means, what a
certified observable family is, what an effective smooth representative
is.  Once those definitions are set correctly, the theorems follow
directly.  This thinness is not a weakness; it reflects that the
book's job is tool-licensing, not physics.  The physics begins in the
branch books, which must come to Book~VI with their observable families
in hand and earn the license.
\end{scholium}

\end{document}
